% Annexe A — Validation des compétences du Bloc 4
\section*{Objet de l'annexe}
Cette annexe reprend, compétence par compétence, la grille du Bloc~4 (RNCP~39583). Pour chacune, elle indique comment la compétence est validée par le maintien en condition opérationnelle de BrackX, en renvoyant à la section du dossier qui la traite, au livrable attendu et aux preuves déposées en annexe. À la date de rédaction (20~août~2026), l'application est en production depuis le 10~août~2026 en version \texttt{v0.6.0}, portée à \texttt{v0.6.1} par le premier correctif.

\subsection*{C4.1.1 --- Mise à jour des dépendances et bibliothèques tierces}
La compétence est validée par un processus de mise à jour outillé, décrit à la section~\ref*{sec:dependances}. La veille est automatisée par Renovate (natif GitLab), qui ouvre des merge requests de mise à jour. Le processus précise sa \textbf{fréquence} (sécurité au fil des avis, mineures hebdomadaires groupées, majeures en revue mensuelle), son \textbf{périmètre} (backend NestJS/Prisma, frontend Angular, image Docker, outillage CI) et le \textbf{type} de traitement (automatique pour les patchs et versions mineures si la CI est verte, manuel avec évaluation d'impact pour les majeures). La sécurité des dépendances est évaluée par Trivy et \texttt{npm audit}. Le livrable — la description du processus — est en section~\ref*{sec:dependances} et ses preuves en \hyperref[chap:annexe-dependances]{annexe~\ref*{chap:annexe-dependances}}.

\subsection*{C4.1.2 --- Système de supervision et d'alerte}
La compétence est validée par un système de supervision adapté à la typologie du logiciel (application web conteneurisée), décrit à la section~\ref*{sec:supervision}. L'outil Rybbit réunit surveillance d'uptime, Core Web Vitals et analytics. Quatre \textbf{sondes} de finalité explicite sont en place (moniteur d'uptime, endpoint \texttt{/health} vérifiant la base via Prisma, collecte des Core Web Vitals, journaux structurés), avec des \textbf{critères de qualité et de performance} chiffrés (uptime, temps de réponse, LCP/INP/CLS, taux d'erreurs) et une \textbf{alerte} sur indisponibilité ou dépassement de seuil. Le système surveille bien la \textbf{disponibilité} du logiciel. Le livrable — la description du système — est en section~\ref*{sec:supervision} et ses preuves en \hyperref[chap:annexe-supervision]{annexe~\ref*{chap:annexe-supervision}}.

\subsection*{C4.2.1 --- Consignation des anomalies}
La compétence est validée par un processus de collecte et de consignation structuré, décrit à la section~\ref*{sec:consignation}. La collecte s'appuie sur quatre canaux (supervision, analyse statique, support, journaux) convergeant vers un point unique~: les issues GitLab, encadrées par un gabarit de bogue et une échelle de sévérité. La \textbf{fiche de consignation} contient les informations permettant de \textbf{reproduire} le bogue, ainsi qu'une analyse de la cause racine et une préconisation de correction. Le livrable — le processus et une fiche de consignation — est illustré par l'anomalie \texttt{\#142}, dont la fiche complète est en \hyperref[chap:annexe-fiche-anomalie]{annexe~\ref*{chap:annexe-fiche-anomalie}}.

\subsection*{C4.2.2 --- Création et déploiement d'un correctif}
La compétence est validée par le traitement de l'anomalie \texttt{\#142}, présenté à la section~\ref*{sec:correctif}. Le correctif \textbf{tire profit du processus d'intégration et de déploiement continu}~: branche dédiée, commit conventionnel référençant l'issue, merge request déclenchant le pipeline (lint, tests dont un test de non-régression, build), puis déploiement par tag sémantique (\texttt{v0.6.1}). Le \textbf{correctif est décrit} (gestion de l'exempt en ronde suisse) et sa résolution est \textbf{vérifiée} en préproduction puis en production via la supervision. Le livrable — la présentation du traitement — est en section~\ref*{sec:correctif} et ses preuves en \hyperref[chap:annexe-correctif]{annexe~\ref*{chap:annexe-correctif}}.

\subsection*{C4.3.1 --- Axes d'amélioration}
La compétence est validée par des recommandations argumentées, présentées à la section~\ref*{sec:amelioration}. Chaque axe découle d'un \textbf{indicateur} (Core Web Vitals et rejeu de session Rybbit, dette technique SonarQube) ou d'un \textbf{retour utilisateur}, et est évalué en termes de \textbf{gain, coût et délai}, puis priorisé selon une logique valeur/effort. Les recommandations sont \textbf{réalistes et réalisables} au regard du projet (aucune refonte) et \textbf{renforcent l'attractivité} du logiciel (performance, fiabilité, engagement). Le livrable est en section~\ref*{sec:amelioration}.

\subsection*{C4.3.2 --- Journal des versions déployées}
La compétence est validée par la tenue d'un journal des versions (\texttt{CHANGELOG.md}), décrit à la section~\ref*{sec:journal-versions}. Aligné sur les tags sémantiques et le format \emph{Keep a Changelog}, il documente pour chaque version ses \textbf{améliorations} (nouvelles fonctionnalités, anomalies corrigées) et ses \textbf{correctifs}, avec renvoi aux issues traitées. Le livrable — un exemplaire du journal — est illustré en section~\ref*{sec:journal-versions} et fourni intégralement en \hyperref[chap:annexe-versions]{annexe~\ref*{chap:annexe-versions}}.

\subsection*{C4.3.3 --- Collaboration avec le support}
La compétence est validée par un exemple de problème résolu en collaboration avec le support, présenté à la section~\ref*{sec:support}. La présentation comporte le \textbf{contexte du retour client} (classement jugé inattendu), la \textbf{résolution apportée} (expertise technique établissant la conformité du comportement, clarification de la documentation et amélioration de transparence) et la \textbf{contribution des parties prenantes} (support, développeur, commanditaire). Le livrable — l'exemple de problème résolu — est en section~\ref*{sec:support} et le fil de l'échange en \hyperref[chap:annexe-support]{annexe~\ref*{chap:annexe-support}}.
